\documentclass{article}
\usepackage[utf8]{inputenc}
\usepackage{amsmath}
\usepackage{booktabs}
\usepackage{multirow}
\usepackage{geometry}
\geometry{a4paper, margin=1in}

\begin{document}

\section*{Airfoil Selection Trade Study}

\subsection*{1. Methodology \& Formulation}
To determine the optimal airfoil among \textbf{E423}, \textbf{GOE165}, and \textbf{S1223} at a Reynolds number $Re = 200,000$, a Weighted Decision Matrix was implemented. Each aerodynamic parameter $i$ is assigned a weight $W_i$ based on mission priorities such that:
\begin{equation}
    \sum_{i=1}^{n} W_i = 1.00
\end{equation}

Each airfoil $j$ receives a discrete score $s_{ij} \in \{-1, 0, 1\}$ representing worst, average/baseline, and best performance, respectively. The total score $S_j$ for airfoil $j$ is calculated as:
\begin{equation}
    S_j = \sum_{i=1}^{n} \left( W_i \cdot s_{ij} \right)
\end{equation}

\subsection*{2. Extracted Aerodynamic Parameters}
The baseline aerodynamic polar data calculated via XFLR5 ($Re = 200,000$) is summarized in Table~\ref{tab:extracted_data}.

\begin{table}[h!]
\centering
\caption{Extracted Aerodynamic Data ($Re = 200,000$)}
\label{tab:extracted_data}
\begin{tabular}{lccccc}
\toprule
\textbf{Airfoil} & \textbf{$C_{L,\max}$} & \textbf{$C_{D,\min}$} & \textbf{$(L/D)_{\max}$} & \textbf{$C_{m0}$ ($\alpha=0^\circ$)} & \textbf{$\alpha_{\text{stall}}$} \\
\midrule
E423   & 2.0338 & 0.01898 & 75.33 & -0.2451 & $13.0^\circ$ \\
GOE165 & 1.5447 & 0.01313 & 60.91 & -0.1587 & $11.5^\circ$ \\
S1223  & 2.2314 & 0.01778 & 71.81 & -0.2716 & $13.0^\circ$ \\
\bottomrule
\end{tabular}
\end{table}

\subsection*{3. Weight Distribution \& Evaluation Criteria}
\begin{itemize}
    \item \textbf{Maximum Lift Coefficient ($C_{L,\max}$, $W_1 = 0.35$):} Critical for payload capacity and take-off performance. Higher is better.
    \item \textbf{Maximum Lift-to-Drag Ratio ($(L/D)_{\max}$, $W_2 = 0.25$):} Defines aerodynamic efficiency and endurance. Higher is better.
    \item \textbf{Minimum Drag Coefficient ($C_{D,\min}$, $W_3 = 0.20$):} Minimizes parasitic drag during dash/cruise. Lower is better.
    \item \textbf{Zero-Lift Pitching Moment ($C_{m0}$, $W_4 = 0.20$):} Affects longitudinal stability and trim drag. Less negative is better.
\end{itemize}

\subsection*{4. Trade Study Decision Matrix}
Table~\ref{tab:trade_matrix} presents the raw ratings, assigned weights, and weighted score calculations.

\begin{table}[h!]
\centering
\caption{Airfoil Trade Study Matrix}
\label{tab:trade_matrix}
\begin{tabular}{lccccccr}
\toprule
\multirow{2}{*}{\textbf{Parameter}} & \multirow{2}{*}{\textbf{Weight ($W_i$)}} & \multicolumn{3}{c}{\textbf{Scores ($s_{ij}$)}} & \multicolumn{3}{c}{\textbf{Weighted Scores ($W_i \cdot s_{ij}$)}} \\
\cmidrule(lr){3-5} \cmidrule(lr){6-8}
& & \textbf{E423} & \textbf{GOE165} & \textbf{S1223} & \textbf{E423} & \textbf{GOE165} & \textbf{S1223} \\
\midrule
$C_{L,\max}$   & 0.35 & 0  & -1 & +1 & 0.00  & -0.35 & 0.35 \\
$C_{D,\min}$   & 0.20 & -1 & +1 & 0  & -0.20 & 0.20  & 0.00 \\
$(L/D)_{\max}$ & 0.25 & +1 & -1 & 0  & 0.25  & -0.25 & 0.00 \\
$C_{m0}$       & 0.20 & 0  & +1 & -1 & 0.00  & 0.20  & -0.20 \\
\midrule
\textbf{Total Score ($S_j$)} & \textbf{1.00} & -- & -- & -- & \textbf{+0.05} & \textbf{-0.20} & \textbf{+0.15} \\
\bottomrule
\end{tabular}
\end{table}

\subsection*{5. Detailed Mathematical Calculations}
The final weighted total scores are computed explicitly as follows:

\begin{align}
    S_{\text{E423}}   &= (0.35 \times 0) + (0.20 \times -1) + (0.25 \times 1) + (0.20 \times 0) = \mathbf{+0.05} \\
    S_{\text{GOE165}} &= (0.35 \times -1) + (0.20 \times 1) + (0.25 \times -1) + (0.20 \times 1) = \mathbf{-0.20} \\
    S_{\text{S1223}}  &= (0.35 \times 1) + (0.20 \times 0) + (0.25 \times 0) + (0.20 \times -1) = \mathbf{+0.15}
\end{align}

\subsection*{6. Selection Conclusion}
According to the trade study results, \textbf{S1223} obtains the highest total score ($S = +0.15$) and is recommended as the optimal choice due to its high lift capability ($C_{L,\max} = 2.2314$).

\end{document}